\chapter{Discussion}
\section{RQ1: Evidence preservation in the implemented workflow}

RQ1 asks how the implemented workflow preserves identity, time, missing states, provenance and
approval boundaries from input to report. The repository and engineering checks show that these
controls exist in one local workflow. Inputs retain dataset and period identifiers; company matches
can enter an explicit hold; missing values keep typed reasons; source records retain locators and
checksums; verifier decisions are persisted; and export requires approval of the current report
version. This is evidence that the workflow is inspectable, not evidence that it is ready for
production \citep{hevner2004design,peffers2007dsrm}.

D0 supplies a narrower behavioural result. Both conditions normalised all fourteen fictional cases
as expected and returned the same ordered outcomes in three repeats. C2 emitted the five supported
cases and held the remaining nine. The accepted observed-zero case shows why a typed zero must not be
collapsed into missing data. The contradicted, stale and untrusted examples show how period, value and
trust rules change the report consequence. These decisions remain traceable to exact fixture records
\citep{gao2023alce,gao2023rarr}.

Being able to trace a claim does not make its source true. A fixed workflow will faithfully repeat a
publisher's mistake, and an exact passage can support a sentence that still misleads. What the
workflow offers is narrower: a reader can see where a value came from, what was done to it, why it
was admitted and who released it. Judging whether the source is right, and how much of the picture it
covers, are separate jobs \citep{gao2023rarr,pineau2021reproducibility}.

RQ1 is therefore answered for the implemented mechanisms under controlled inputs: the workflow
produces a complete record of why each claim was emitted or withheld, and Section~7.5 states the
boundary of that answer \citep{pineau2021reproducibility,gebru2021datasheets,nist2023airmf}.

\section{RQ2: Observed effect of separate verification}

RQ2 compares C1 and C2 on identical D0 candidates and evidence. C1 emitted eleven claims, including
six designed to be unsupported. Its precision was 0.455 and its unsupported-claim rate was 0.545. C2
emitted five claims, all labelled as supported, giving precision of 1.000 and an unsupported-claim
rate of 0.000. Recall remained 1.000 in both conditions. The observed effect was selective admission,
not discovery of additional evidence \citep{gao2023alce,gao2023rarr}.

The lower claim count is important. C2 did not improve by producing a more complete report; it
withheld candidates for conflicting values, a stale period, absent or inaccessible support, mixed
currency and untrusted content. This was correct for D0, but a similar gate could falsely reject a
valid claim in new data. D0 cannot estimate false rejection outside its visible designed cases
\citep{kapoor2023leakage,nist2023airmf}.

The scope of that result needs stating precisely, and it is exactly the scope
Section~\ref{sec:design-alternatives} claimed in advance. C1 and C2 share one program, one catalogue
and one normalisation route, so what differs between them is the support rule and its authority over
composition, not a persona, a conversation style or a count of agents. A deterministic pipeline or a
single grounded agent could apply the same rule and reach the same decisions. The finding is that an
explicit, separately recorded refusal changed what the report was allowed to say; it is not evidence
that multi-agent designs are generally better \citep{guo2024multiagent,wu2024autogen}.

No useful speed or cost comparison was produced. Both routes used local fixed rules, model cost was
zero and sub-millisecond times were truncated to zero. Role separation also adds states, hand-offs,
storage, tests and maintenance. A later pilot must measure these costs alongside false acceptance,
false rejection and reviewer work \citep{pineau2021reproducibility,hevner2004design}.

\section{Why functional role separation was selected}

The selected design separates production from verification and keeps approval outside both. The
planner and collectors can propose work, but they cannot approve a report. The verifier can accept,
hold or contradict a candidate, but it cannot rewrite the claim to make it pass. The composer can use
only admitted claims, and the reviewer decides whether the current version may be exported. Persisted
hand-offs make each decision recoverable and allow a failure to stop one stage without silently
changing another \citep{guo2024multiagent,wu2024autogen}.

This design also supports narrower permissions. Public discovery, protected collection, local
restricted-data processing and final approval do not share one unrestricted tool boundary. The trade-
off is more code and more operational state. Multi-agent trace studies report coordination and task-
verification failures, so added roles are themselves a source of risk
\citep{cemri2025masfailures,nist2023airmf}. The prototype uses a fixed state machine rather than an
open-ended conversation to keep this cost visible.

The alternatives in Table~\ref{tab:lit-solution-landscape} remain credible, for the reasons
Section~\ref{sec:design-alternatives} set out. This workflow earns its place by making the support
and approval boundaries explicit and open to inspection, not by being the best option for every
reporting team \citep{gao2023ragsurvey,guo2024multiagent}.

\section{Public-web case study and future comparator}

The public-web route extends the implementation from structured replay to bounded source discovery:
web search proposes addresses, local controls decide whether a page may be collected, and a claim
still needs a saved source, an exact supporting passage and approval
\citep{openai2026websearch,gao2023alce}. Controlled tests covered reviewed identity, network checks,
request limits, write-once captures, redaction, exact passages, contradictions and approval using fake
providers and fixtures, which establishes behaviour for those cases and nothing about the open web
\citep{greshake2023indirect,nist2023airmf}.

The most honest challenge to the whole custom build is that a managed platform might do it for less
work. Grok Bot's documentation describes persistent named Bots with defined jobs, tools, hand-offs and
approval boundaries, covering a good deal of what Chapter~4 implements by hand
\citep{spacexai2026grokbot}. Lower effort remains a hypothesis, and the decisive question for any such
comparison is whether the platform can refuse a claim the evidence does not support, because a
configurable Bot is not by itself a verification contract
\citep{dibia2024autogenstudio,peffers2007dsrm}.

\section{Relationship to literature and design trade-offs}

The findings agree with literature that separates factual correctness, citation support and source
attribution. Exact locators and saved copies make a decision inspectable, but they do not establish
source truth or complete coverage \citep{gao2023alce,gao2023rarr}. The observed D0 change is consistent
with the purpose of a verification gate because the fixture was built to test that gate. It is not an
independent confirmation of a general theory.

Fixed missing-state, cut-off and version rules also support repeatability. Their weakness is that a
wrong or outdated rule repeats consistently. The rules therefore need named ownership, change
control and re-evaluation when sources, schemas or reporting needs change
\citep{pineau2021reproducibility,nikiforova2020quality,gebru2021datasheets}.

Human approval provides accountability for release, not automatic correctness. Interfaces should
show uncertainty and allow correction, while studies of reliance show that people may still follow
incorrect advice \citep{amershi2019guidelines,bucinca2021forcing}. The pilot must measure reviewer
edits, decisions and time instead of treating an approval button as evidence of benefit.
Table~\ref{tab:disc-rq-evidence-status} in Appendix~\ref{app:supplementary-evaluation} separates the
two answered research questions from the unexecuted comparisons. Table~\ref{tab:disc-literature-result-alignment}
in Appendix~\ref{app:supplementary-exhibits} shows where the results support, qualify or leave the
literature claims untested.

\section{Generalisability and practical interpretation}

Everything above applies to one local prototype, fourteen designed cases and a set of controlled
security tests. Early-stage private-market evidence is unusually incomplete and sensitive to how terms
are defined, which is precisely why statistical generalisation would need an independently selected
and labelled set rather than a larger version of D0
\citep{kaplan2016vcdata,britishbusinessbank2025equity,kapoor2023leakage}.

Moving the design to another setting is not a matter of copying code. It would need a new data
dictionary, a legal identity source, source licences, missing-value rules, a reporting cut-off, an
evidence contract, a security boundary and a named approval owner. Checksums and versions help explain
why two runs differ; they do not hold behaviour steady when the sources or models underneath change
\citep{gebru2021datasheets,nist2023airmf}. The transfer checklist is
Table~\ref{tab:disc-transfer-conditions} in Appendix~\ref{app:supplementary-exhibits}.

The prototype is best understood, then, as a reporting control system that can be tested. It ranks
nothing, predicts nothing and advises nobody, and whether it is genuinely useful remains a question
for the pilot \citep{fca2024promotions,hevner2004design}.

\section{Prospective business pilot}

The proposed pilot is a time-limited evaluation of the existing portfolio XLSX import as the named
input integration. The prototype's Next.js review interface is the downstream report interface, with
versioned JSON and accessible HTML exported only after approval. The pilot owner approves scope and
the go/no-go decision. A reporting analyst prepares the manual baseline and operates each reporting
case. A technical operator manages the local service, versions and failures. A reviewer or approver
checks the evidence and current report version. A data and security owner approves data classes,
source licences, storage, retention and any external-model route
\citep{amershi2019guidelines,nist2023airmf}.

Before data collection, the team freezes the pilot's size as well as its rules. Four figures fix the
size: how many portfolio companies, how many reporting periods, what calendar window, and the
smallest number of completed paired reports that would make the comparison worth reading. The pilot
owner and the reviewer choose those four; this dissertation leaves them open, because a size picked
without knowing the site's caseload would be invented rather than planned. The team also freezes the
eligible sources, schema, reference decisions, task completion rule and acceptance thresholds. It
records setup effort in
staff-hours for data mapping, source admission, environment configuration, training and dry runs.
For every report it records analyst active minutes, reviewer active minutes, elapsed time, review and
rework time, retries, model or API tokens, provider charges, source-licensing charges, storage and
security costs. Cost per report is the allocated setup cost plus staff time, API use, licences,
storage and security, divided by completed reports. No saving is assumed before this matched baseline
exists \citep{pineau2021reproducibility,peffers2007dsrm}.

Quality measures are supported-claim precision, evidence coverage, false acceptance, false rejection,
reviewer edits, completion rate and source failures. Operational measures are active time, elapsed
time and cost per report. The manual baseline uses the same inputs, source window and completion rule.
Results are paired by company-period, and failed or abandoned cases remain in the denominator where
the measure requires them. Data-quality measures must remain tied to the intended reporting use
\citep{nikiforova2020quality}, and reviewer disagreement must not be treated as proof that either
person is correct \citep{artstein2008agreement}.

The pilot proceeds only if authority, source licences, reference labels, storage controls, trained
owners and rollback are in place. Before seeing outcomes, the pilot owner and reviewer set the
minimum acceptable quality and completion values and the maximum acceptable time, false-acceptance,
false-rejection and cost values. Time is on that list because a saving nobody committed to in advance
cannot later be claimed as one. The recommendation is go only when those frozen criteria are met,
approval cannot be bypassed, every exported claim has the required evidence record, and no stop event
occurs. It is no-go or redesign when a criterion fails. This avoids inventing numeric targets in the
dissertation while leaving a checkable decision rule
\citep{peffers2007dsrm,nist2023airmf}.

The run stops immediately for unclear authority or consent, unexpected restricted or personal data,
credential exposure, a source-licence or robots breach, external transmission outside the approved
route, a broken evaluation freeze, evidence-record loss, approval bypass, uncontrolled cost or a
security incident. The event, affected cases and recovery decision are recorded. Expected time or
cost reduction remains a hypothesis until the paired pilot demonstrates it
\citep{cddo2023genai,nist2023airmf}.

Four kinds of testing remain outstanding, and they are gathered here so that a pilot owner can find
them in one place. An independent tester needs to attack the collection and extraction route with
hostile inputs. An accessibility assessment needs assistive technology and an assessor from outside
the project. A security assessment needs to cover remote and multi-user operation, neither of which
this prototype supports. And a live-source study needs to run the approved official connectors
against real registry and publisher data. Sections~7.5 and~8.3 refer to this list rather than
repeating it \citep{nist2023airmf,pineau2021reproducibility}.
